\documentclass[11pt]{article}
\usepackage[margin=1in]{geometry}
\usepackage{amsmath,amssymb,amsthm,mathtools}
\usepackage{graphicx}
\usepackage{xcolor}
\usepackage{enumitem}
\usepackage[hidelinks]{hyperref}
\usepackage{microtype}

\newtheorem{theorem}{Theorem}
\newtheorem{lemma}[theorem]{Lemma}
\newtheorem{proposition}[theorem]{Proposition}
\newtheorem{remark}[theorem]{Remark}
\newcommand{\K}{\mathbb K}
\newcommand{\R}{\mathbb R}
\newcommand{\C}{\mathbb C}
\newcommand{\cA}{\mathcal A}
\newcommand{\ip}[2]{\langle #1,#2\rangle}

\title{Expander-Sparsified Tensor Embeddings Give a\\
Deterministic Rank-One Matrix RIP}
\author{Discovery packet for arXiv:2310.18565\\
\texttt{fa\_banach\_001}, lane 17, \texttt{agent\_lane\_17}}
\date{11 August 2026}

\begin{document}
\maketitle

\begin{center}
\fbox{\parbox{0.92\textwidth}{
\textbf{Status.} Candidate substantial partial result, likely valid, requiring
expert review.  It answers the source challenge for matrices of rank one, over
both $\R$ and $\C$, with rank-one measurement matrices and
$m=n^{1+o(1)}$.  The case of rank $r>1$ remains open.
}}
\end{center}

\section{Source and exact challenge}

The source is Simon Foucart, \emph{Linearly Embedding Sparse Vectors from
$\ell_2$ to $\ell_1$ via Deterministic Dimension-Reducing Maps},
arXiv:2310.18565, Section~6, page~14 \cite{Foucart}.  It asks for a
deterministic linear map
\[
  \cA:\K^{n\times n}\longrightarrow \K^m
\]
such that, for all matrices $X$ of rank at most $r$,
\begin{equation}\label{eq:source-rip}
 (1-\delta)\|X\|_F\leq \|\cA(X)\|_1
 \leq(1+\delta)\|X\|_F,
\end{equation}
in a regime $m\asymp n^\lambda r^\mu\ll n^2$, preferably with rank-one
measurement matrices.  The source calls this a ``wide open question.''

\begin{figure}[p]
  \centering
  \includegraphics[width=\textwidth]{figures/open_problem_crop.png}
  \caption{The exact challenge in arXiv:2310.18565, page 14.  The image is
  rendered from the local source PDF included with this packet.}
\end{figure}
\clearpage

We use the Frobenius pairing $\ip{A}{X}_F=\operatorname{tr}(A^*X)$.
An $M\times n$ matrix $B$, whose $j$th row is $b_j^*$, will be called
\emph{$(a,b)$-flat and tight} if
\begin{equation}\label{eq:flat}
 B^*B=I_n,
 \qquad
 a\sqrt M\|x\|_2\leq \|Bx\|_1
 \leq b\sqrt M\|x\|_2
 \quad(x\in\K^n).
\end{equation}

\section{Result}

\begin{theorem}[Rank-one subcase of the modified matrix RIP]\label{thm:main}
Let $\K\in\{\R,\C\}$ and fix $0<\delta<1$.  There is a deterministic
polynomial-time construction, for every sufficiently large $n$, of rank-one
matrices $A_1,\dots,A_m\in\K^{n\times n}$ such that
\begin{equation}\label{eq:main-rip}
 (1-\delta)\|X\|_F
 \leq \sum_{i=1}^m |\ip{A_i}{X}_F|
 \leq (1+\delta)\|X\|_F
\end{equation}
for every matrix $X\in\K^{n\times n}$ of rank at most one.  Moreover,
\begin{equation}\label{eq:m-bound}
 m\leq n\exp\!\bigl(C_\delta(\log\log(n+2))^2\bigr)
      =n^{1+o(1)}=o(n^2).
\end{equation}
The finitely many smaller dimensions can be covered deterministically after
enlarging the constant depending on $\delta$.
\end{theorem}

\section{Proof intuition}

An explicit almost-Euclidean subspace of $\ell_1^M$ supplies a matrix $B$
whose vector $Bx$ has almost constant $\ell_1$-to-$\ell_2$ ratio.  The
specific recursive extractor construction of Indyk has an additional feature:
each stage repeats coordinates of a concatenation of orthonormal bases with
constant multiplicity.  Hence it is an exact tight frame.  After normalization,
$B$ obeys \eqref{eq:flat} with $b/a$ as close to one as desired and
$M=n^{1+o(1)}$.

The full tensor coordinates of a rank-one matrix $X=uv^*$ factor:
\[
  |\ip{b_jb_k^*}{uv^*}_F|
       =|b_j^*u|\,|b_k^*v|.
\]
There is therefore no need to keep all $M^2$ pairs $(j,k)$.  Retain only the
edges of a constant-degree spectral expander on the two coordinate sets.
The expander mixing lemma compares the retained sum with the product of the
two $\ell_1$ norms, and tightness controls the error by
$\|u\|_2\|v\|_2=\|uv^*\|_F$.  Choosing the frame distortion and normalized
second eigenvalue small enough gives any prescribed $\delta$.  For rank
$r>1$, absolute values destroy this separability; that is exactly where the
argument stops.

\section{Two ingredients}

\begin{lemma}[Explicit flat tight frames]\label{lem:flat}
For every $0<\eta<1/4$ and $\K\in\{\R,\C\}$ there are constants
$0<\rho_0\leq\rho_1<\infty$, independent of $n$ and $\eta$, and a
deterministic polynomial-time construction of $B=B_{n,\eta}\in\K^{M\times n}$
such that
\begin{align}
 B^*B&=I_n,\label{eq:tight}\\
 (1-\eta)\rho\sqrt M\|x\|_2
 &\leq \|Bx\|_1
 \leq(1+\eta)\rho\sqrt M\|x\|_2,\label{eq:nearly-flat}
\end{align}
for some $\rho\in[\rho_0,\rho_1]$, and
\begin{equation}\label{eq:flat-dim}
 M\leq n\exp\!\bigl(C_\eta(\log\log(n+2))^2\bigr).
\end{equation}
\end{lemma}

\begin{proof}
We record why the standard explicit embedding can be chosen tight; this is
the structural point used later.  Indyk's Theorem~1.1 and Corollary~5.2
construct an explicit $\ell_2^n\to\ell_1^M$ embedding with distortion
$1+O(\eta)$ and $M=n\exp(O_\eta((\log\log n)^2))$ by recursively applying
extractor-based block maps \cite{Indyk}.  At one recursion level its dictionary
$D$ is the vertical concatenation of $L$ orthonormal matrices, so
$D^*D=LI$.  The extractor graph is left $d$-regular, and the associated map
$F$ merely places every coordinate of $Dx$ in exactly $d$ output positions.
Consequently
\begin{equation}\label{eq:stage-tight}
 F^*F=LdI.
\end{equation}
After division by $\sqrt{Ld}$, each recursion level is an exact isometry.
Indyk's displayed block estimate becomes
\begin{equation}\label{eq:block}
 (1-O(\zeta))\sqrt b\,\|x\|_2
 \leq\sum_{s=1}^b\|(Fx)_s\|_2
 \leq\sqrt b\,\|x\|_2
\end{equation}
for the normalized map, where the output is split into $b$ blocks.

Iterate \eqref{eq:block} on every block for $O(\log\log n)$ levels, taking
$\zeta$ a sufficiently small multiple of $\eta/\log\log(n+2)$.  The product
of the lower factors is at least $1-\eta/2$, while every level remains an
exact isometry by \eqref{eq:stage-tight}.  At the terminal fixed block
dimension, use a deterministic finite spherical quadrature.  Concretely, a
sufficiently fine finite set on the sphere can be symmetrized under signed
coordinate permutations.  Its covariance is then exactly scalar, while its
equal-weight averages of the equicontinuous family
$y\mapsto|\ip{x}{y}|$ approximate the rotational average uniformly in
$x$.  Repeating points clears rational weights.  After scalar normalization,
this terminal map is tight and has $\ell_1$ distortion at most
$1+\eta/2$.  Its dimension depends only on $\eta$.

Composing the normalized block maps with the terminal maps gives
\eqref{eq:tight}.  The number of leaves is $M$; multiplying the $\sqrt b$
factors in \eqref{eq:block} gives exactly $\sqrt M$ (up to the fixed terminal
factor).  This proves \eqref{eq:nearly-flat}.  The rotational first absolute
moment times the square root of the terminal dimension is bounded above and
below by universal positive constants, giving $\rho_0,\rho_1$.  Indyk's
dimension calculation gives \eqref{eq:flat-dim}.  Padding the input to the
next admissible dictionary dimension and restricting the columns preserves
tightness and changes only the displayed constant.

The construction above is real.  To retain the same distortion over $\C$,
apply it in real dimension $2n$ and write its row matrix as $R=[P\ Q]$.
For $z=x+iy$, put
\[
 C=P-iQ,\qquad B_\C=2^{-1/2}C.
\]
Since $R^TR=I_{2n}$, we have
$P^TP=Q^TQ=I_n$ and $P^TQ=Q^TP=0$, hence $C^*C=2I_n$ and
$B_\C^*B_\C=I_n$.  If $U_\theta$ is realification of multiplication by
$e^{i\theta}$, then, up to the harmless choice of sign of $\theta$,
\[
 (RU_\theta z)_j=\operatorname{Re}(e^{i\theta}(Cz)_j).
\]
Integrating the real flat estimate in $\theta$ and using
\[
 \frac1{2\pi}\int_0^{2\pi}|\operatorname{Re}(e^{i\theta}w)|\,d\theta
   =\frac{2}{\pi}|w|
\]
shows that $B_\C$ has exactly the same relative $1+O(\eta)$ distortion;
only $\rho$ is multiplied by $\pi/(2\sqrt2)$.  The dimension is
$(2n)^{1+o(1)}$, proving the complex statement as well.
\end{proof}

\begin{lemma}[All-orders explicit expander mixing]\label{lem:expander}
For every $\theta>0$ there is a constant $d=d(\theta)$ such that, for every
sufficiently large $M$, one can deterministically construct in time
polynomial in $M$ a $d$-regular bipartite graph $G=(L,R,E)$ with
$|L|=|R|=M$ whose adjacency matrix has all singular values except the top one
at most $\theta d$.  Thus, for all $f,g\in\R^M$,
\begin{equation}\label{eq:mixing}
 \left|\sum_{(j,k)\in E}f_jg_k
 -\frac dM\left(\sum_jf_j\right)\left(\sum_kg_k\right)\right|
 \leq \theta d\|f\|_2\|g\|_2.
\end{equation}
\end{lemma}

\begin{proof}
Alon gives a deterministic polynomial-time construction, on exactly every
sufficiently large number $M$ of vertices, of a $d$-regular graph with every
nontrivial eigenvalue at most $2\sqrt{d-1}+o_d(1)$ in absolute value
\cite{Alon}.  First choose the constant $d$ large enough and then the finite
threshold for $M$ so that this is at most $\theta d$.  Its bipartite double
cover has the same adjacency matrix between its two vertex classes, hence the
required singular-value bound.  Decomposing $f$ and $g$ into their constant
and mean-zero parts proves \eqref{eq:mixing}.
\end{proof}

\section{Proof of Theorem~\ref{thm:main}}

Fix a small $\eta>0$, to be chosen from $\delta$, and take the frame in
Lemma~\ref{lem:flat}.  Set
\[
 a=(1-\eta)\rho,\qquad b=(1+\eta)\rho.
\]
Choose $0<\theta<a^2$, and let $G$ be the graph in
Lemma~\ref{lem:expander}.  For every edge $e=(j,k)$ define the rank-one
measurement matrix
\begin{equation}\label{eq:measurement}
 A_e=b_jb_k^*.
\end{equation}

Let $X=uv^*$ have rank at most one.  (The zero case is immediate.)  With
$f_j=|b_j^*u|$ and $g_k=|b_k^*v|$, the Frobenius pairing factors as
\begin{equation}\label{eq:factor}
 |\ip{A_e}{X}_F|=|b_j^*u|\,|b_k^*v|=f_jg_k.
\end{equation}
Tightness gives $\|f\|_2=\|u\|_2$ and $\|g\|_2=\|v\|_2$, while flatness
gives
\[
 a\sqrt M\|u\|_2\leq\|f\|_1\leq b\sqrt M\|u\|_2,
 \qquad
 a\sqrt M\|v\|_2\leq\|g\|_1\leq b\sqrt M\|v\|_2.
\]
Apply \eqref{eq:mixing} to the nonnegative vectors $f,g$.  Since
$\|uv^*\|_F=\|u\|_2\|v\|_2$, it follows that
\begin{equation}\label{eq:raw-bounds}
 d(a^2-\theta)\|X\|_F
 \leq \sum_{e\in E}|\ip{A_e}{X}_F|
 \leq d(b^2+\theta)\|X\|_F.
\end{equation}
Put $\alpha=d(a^2-\theta)$ and $\beta=d(b^2+\theta)$.  Multiplying every
$A_e$ by $2/(\alpha+\beta)$ turns \eqref{eq:raw-bounds} into
\eqref{eq:main-rip} with
\begin{equation}\label{eq:delta-prime}
 \delta'=\frac{\beta-\alpha}{\beta+\alpha}
 =\frac{b^2-a^2+2\theta}{a^2+b^2}.
\end{equation}
First choose $\eta$ and then $\theta$ sufficiently small that
$\delta'\leq\delta$.  The number of edges is
$m=dM$, and $d$ depends only on $\delta$.  Lemma~\ref{lem:flat} gives
\eqref{eq:m-bound}.  Every matrix in \eqref{eq:measurement} is rank one and
all ingredients are deterministic and polynomial-time constructible.
\qed

\section{Why the proof does not yet extend beyond rank one}

Six focused upgrade routes were checked after obtaining the rank-one theorem.
They delimit the current result rather than form part of its proof.

\begin{enumerate}[leftmargin=*,label=\textbf{Attempt \arabic*.}]
\item \textbf{SVD plus triangle inequalities.}
For $X=\sum_{s=1}^r\sigma_su_sv_s^*$, the upper estimate follows with
$\|X\|_*$ and hence loses at most $\sqrt r$.  The lower estimate fails because
the summands cancel inside each scalar measurement before the absolute value
is taken.

\item \textbf{Dense tensor coordinates first.}
Keeping all $M^2$ entries of $BXB^*$ yields a deterministic embedding at the
trivial quadratic scale.  One can prove a constant lower bound by applying
the flat estimate to columns and averaging scalar projections.  This does
not reduce the measurement count below $n^{2+o(1)}$.

\item \textbf{Ordinary spectral sampling of $|BXB^*|$.}
For rank one, this nonnegative matrix is an outer product.  For rank $r>1$,
entrywise absolute value can have full rank.  The generic trace-norm estimate
in the mixing error loses $\sqrt M$; compensating for it forces degree of
order $M$ and returns to $M^2$ measurements.

\item \textbf{Positive semidefinite restriction.}
For $X\succeq0$, tight-frame quadratic measurements sum to
$\operatorname{tr}X$.  This only gives
$\|X\|_F\leq\operatorname{tr}X\leq\sqrt r\|X\|_F$, not a prescribed small
RIP constant, and it does not treat indefinite or nonsymmetric matrices.

\item \textbf{Phase lifting and sign-rank samplers.}
The identity $|z|=\sup_{|\omega|=1}\operatorname{Re}(\omega z)$ introduces
an independent phase at every sampled edge.  Those edgewise phases destroy
the rank-$r$ factorization.  A sampler uniform over this phase-enlarged class
would be a genuinely stronger pseudorandom object than a spectral expander.

\item \textbf{Polynomial/moment approximation to the absolute value.}
Approximating $|z|$ by polynomials in $|z|^2$ turns the problem into sampling
higher tensor powers.  Without a uniform entry bound, the approximation is
not relative near zero; after truncation, the known tensor designs have
degree or size exponential in the approximation degree or in $r$.  No bound
of the requested form $n^\lambda r^\mu$ emerges.
\end{enumerate}

Thus the surviving obstruction is precise: a graph expander samples
separable bilinear arrays, whereas $|BXB^*|$ is no longer separable after the
rank-one case.  A low-rank absolute-value sampler with polynomial dependence
on $r$ would upgrade the theorem, but no such lemma has been established here.

\section{Verification, novelty bounds, and review recommendation}

The accompanying script \texttt{code/verify\_rank\_one\_expander.py} builds a
finite exact flat frame from all sign vectors, a deterministic union-of-
permutations bipartite graph, and the rank-one measurements
\eqref{eq:measurement}.  It checks tightness, the second singular value, the
factorization identity, the raw constants in \eqref{eq:raw-bounds}, and the
rescaled RIP bounds on a fixed reproducible suite of real and complex samples.
This is a regression check, not a proof; the proof is the two-lemma argument
above.

The four lightweight run indexes were searched for arXiv:2310.18565, the
title, ``modified RIP,'' ``rank-one measurements,'' deterministic low-rank
recovery, tensor embeddings, and expander mixing.  They contained adjacent
records about deterministic sparse-vector RIP (arXiv:1003.2990, 1202.1234,
and 1406.4089) but no statement of the theorem above.  A bounded primary-source
search of arXiv and citation/phrase searches through 11 August 2026 located
the source, Indyk's embedding, standard explicit expanders, and randomized
rank-one matrix RIP work, but no later paper stating this expander-sparsified
rank-one construction.  This is not an exhaustive novelty certification.

The main human-review focus should be Lemma~\ref{lem:flat}: verify that the
tight normalization is retained through the recursive block construction and
terminal fixed-dimensional quadrature.  The expander step is then a direct
one-line application of mixing.  The claimed status should remain
\emph{candidate substantial partial result} until that structural extraction
and the bounded novelty search receive independent expert review.

\begin{thebibliography}{9}

\bibitem{Foucart}
S. Foucart,
\emph{Linearly Embedding Sparse Vectors from $\ell_2$ to $\ell_1$ via
Deterministic Dimension-Reducing Maps},
arXiv:2310.18565 (2023), Section 6.

\bibitem{Indyk}
P. Indyk,
\emph{Uncertainty Principles, Extractors, and Explicit Embeddings of $L_2$
into $L_1$},
Electronic Colloquium on Computational Complexity, Report TR06-126 (2006);
also STOC 2007.  See Theorem 1.1 and Corollary 5.2.

\bibitem{Alon}
N. Alon,
\emph{Explicit Expanders of Every Degree and Size},
Combinatorica 41 (2021), 447--463; arXiv:2003.11673.

\end{thebibliography}

\end{document}
